\chapter[1923 --- Banting et al.]{1923: Frederick Banting, John Macleod}
\label{ch:1923_Banting}

\section*{Main Contribution}

\textit{Discovered insulin, transforming diabetes from a fatal disease into a manageable condition---one of the most important medical breakthroughs of the twentieth century.}

\section{Official Citation}

for their discovery of insulin

\section{Background and Historical Context}

\subsection{Frederick Grant Banting}

Frederick Grant Banting was born on November 14, 1891, in Alliston, Ontario, Canada, the youngest of five children in a farming family. Banting attended the University of Toronto, where he initially studied arts before transferring to medicine. He graduated in 1916 with a Bachelor of Medicine degree and served as a medical officer in the Canadian Army Medical Corps during World War I, seeing action in France and Belgium. His wartime experiences, particularly his exposure to the suffering of wounded soldiers, deepened his commitment to medical research. After the war, Banting returned to Canada and began a career in general practice in London, Ontario, while simultaneously pursuing research in physiology at the University of Toronto under the mentorship of Professor John James Rickard Macleod. It was during this period that Banting conceived the idea that would lead to the discovery of insulin---an idea that would save millions of lives and transform the treatment of diabetes.

Banting was not a biochemist or a physiologist by training; he was a general practitioner with a keen clinical mind and an ability to see connections that more specialized scientists had overlooked. His background in general medicine gave him a perspective that was different from that of the laboratory researchers who had been studying diabetes, and it was precisely this outsider's perspective that enabled him to formulate the critical insight that led to the discovery of insulin. Banting's story is a powerful reminder that breakthrough discoveries in medicine sometimes come not from the most experienced or specialized researchers but from individuals who bring fresh eyes and creative thinking to long-standing problems.

\subsection{John James Rickard Macleod}

John James Rickard Macleod was born on September 6, 1876, in Cluny, Scotland. He studied medicine at the University of Aberdeen, where he developed a strong interest in carbohydrate metabolism. After holding positions at the London Hospital Medical College and Western Reserve University in Cleveland, Ohio, Macleod was appointed professor of physiology at the University of Toronto in 1918. By the time Banting arrived in his laboratory, Macleod was already an internationally recognized expert on diabetes and carbohydrate metabolism, having published extensively on the subject and having developed experimental techniques for studying the pancreas and blood sugar regulation. Macleod's laboratory at Toronto was one of the best-equipped in North America, with facilities for physiological and biochemical research that were notable at most other institutions.

Macleod was a meticulous scientist and a demanding supervisor who insisted on rigorous experimental design and careful interpretation of results. His expertise in carbohydrate metabolism and his experience with pancreatic physiology made him an ideal collaborator for Banting, whose clinical insights needed the support of a well-equipped laboratory and a scientifically rigorous environment to be translated into experimental results.

The discovery of insulin took place against the backdrop of a global diabetes epidemic. Before insulin, a diagnosis of type 1 diabetes was essentially a death sentence. The disease, which typically affected young people, caused progressive wasting, ketoacidosis, coma, and death within months of diagnosis. The only known treatment was the so-called ``starvation diet,'' devised by the New York physician Frederick Allen, which restricted caloric intake to the bare minimum necessary to keep blood glucose levels low enough to prevent fatal ketoacidosis. Patients on this regimen survived for months or occasionally a few years, but they were slowly wasting away, and the quality of life was dismal. The wards of hospitals around the world were filled with emaciated children and young adults, their bodies wasting from the inability to utilize the glucose they consumed, their breath smelling of acetone from the ketone bodies that accumulated in their blood.

The search for the pancreatic hormone that controlled blood sugar had been ongoing for decades. In 1889, the German physiologists Josef von Mering and Oskar Minkowski had shown that surgical removal of the pancreas in dogs caused a condition resembling diabetes, proving that the pancreas played an essential role in blood sugar regulation. In the early 1900s, several investigators---including the Romanian physician Nicolas Paulescu and the American scientists Isra{\"e}l Kleiner and John R. Murlin---had attempted to extract the active principle from pancreatic tissue, but all had failed to produce a preparation that was effective and non-toxic when injected into diabetic animals. The pancreas contains two functionally distinct tissue types: the exocrine tissue, which secretes digestive enzymes into the gut, and the endocrine tissue (the islets of Langerhans), which secretes hormones directly into the bloodstream. The problem was that the digestive enzymes destroyed the hormone during the extraction process, rendering all previous attempts at isolation futile.

\section{The Discovery and Contribution}

The story of the discovery of insulin is one of the most dramatic in the history of medicine. In the spring of 1920, Banting read an article in \emph{Surgery, Gynecology and Obstetrics} by the American surgeon Moses Barron, who described a case in which a gallstone had obstructed the pancreatic duct, causing atrophy of the exocrine tissue while leaving the islets of Langerhans relatively intact. This observation gave Banting an idea: if the exocrine tissue were destroyed, the digestive enzymes it produced would no longer interfere with the extraction of the insulin from the islets. By ligating the pancreatic ducts of dogs and waiting for the exocrine tissue to atrophy, one might be able to extract the blood-sugar-regulating substance from the remaining islet tissue.

On the night of July 30, 1920, Banting jotted down his ideas in a notebook, and the next morning he presented them to Macleod. Macleod was initially skeptical---he had heard many such proposals before and had himself attempted to extract the pancreatic hormone without success---but he agreed to provide Banting with laboratory space, experimental animals, and the assistance of a graduate student, Charles Best. The collaboration between Banting and Best, which began on May 17, 1921, was one of the most productive partnerships in the history of medicine.

The early experiments were fraught with difficulties. Banting and Best ligated the pancreatic ducts of a series of dogs and waited for the exocrine tissue to degenerate (a period of approximately eight weeks). They then extracted the atrophied pancreatic tissue and injected the resulting fluid into dogs that had been rendered diabetic by removal of their pancreata. The initial results were promising but inconsistent: some dogs showed a significant drop in blood glucose levels, while others showed only a transient or negligible response. Banting was discouraged by the variability of the results, but Best encouraged him to continue, and they refined their extraction technique, improving the potency of the extract.

The breakthrough came in the fall of 1921, when Banting and Best obtained consistent and dramatic results. Diabetic dogs that received injections of the pancreatic extract showed a rapid and sustained decrease in blood glucose levels, and the symptoms of diabetes---including glycosuria (glucose in the urine), polyuria (excessive urination), and polydipsia (excessive thirst)---were reversed. The most dramatic demonstration of the effectiveness of the extract came on December 20, 1921, when Banting and Best injected a diabetic dog named ``Marjorie'' with the extract and observed a drop in blood glucose from 0.46\% to 0.18\% within three hours---a result that left no doubt about the efficacy of the extract.

In January 1922, the first clinical trial of insulin was conducted at the Toronto General Hospital. The patient was Leonard Thompson, a 14-year-old boy with severe type 1 diabetes who was near death. The initial injection of the extract, prepared by James Collip (a biochemist who had joined the team to help purify the extract), caused an allergic reaction and had only a transient effect. However, after Collip refined the purification procedure, subsequent injections produced sustained improvement in the boy's condition. Leonard Thompson's recovery was dramatic: within days, his blood glucose levels had fallen to near-normal, his appetite returned, and he was able to leave the hospital. He would go on to live another 13 years, dying at the age of 27 from pneumonia---a remarkable outcome for a patient who, without insulin, would have died within months.

The news of the Toronto discovery spread rapidly through the medical world. Within months, laboratories around the world were producing insulin, and by 1923, the pharmaceutical company Eli Lilly and Company had developed methods for large-scale production of the hormone from porcine and bovine pancreata. The impact on patients was notable: children and young adults who had been wasting away from diabetes were able to resume normal lives, and the death rate from diabetes in the United States dropped by more than 50 percent within the first decade of insulin's availability.

Banting and Macleod were awarded the Nobel Prize in Physiology or Medicine in 1923 for the discovery of insulin. The award was controversial: Best and Collip, who had played essential roles in the discovery and purification of insulin, were not included, and there was considerable acrimony between Banting and Macleod over the distribution of credit. Banting was so aggrieved by the omission of Best that he announced he would share his prize money with Best, and Macleod did the same with Collip. Despite these personal disputes, the scientific achievement was universally recognized as a notable medical breakthroughs of the twentieth century.

\section{Impact on Modern Medicine}

The discovery of insulin transformed diabetes from a fatal disease into a manageable chronic condition. Today, approximately 537 million adults worldwide live with diabetes, and insulin remains the cornerstone of treatment for type 1 diabetes and many cases of type 2 diabetes. The development of insulin therapy has progressed far beyond the crude pancreatic extracts of the 1920s. Modern insulin preparations include rapid-acting analogs (lispro, aspart, glulisine), which begin to work within 10--15 minutes and peak within 1--2 hours; long-acting analogs (glargine, detemir), which provide a steady level of insulin over 24 hours; and ultra-long-acting analogs (degludec), which have a duration of action exceeding 42 hours. These preparations allow for precise control of blood glucose levels with reduced risk of hypoglycemia, enabling patients with diabetes to achieve near-normal glycemic control and to live full, active lives.

The insulin pump, first developed in the 1970s and now in widespread use, delivers continuous subcutaneous insulin infusion that mimics the physiological pattern of insulin secretion from the pancreatic beta cells. Closed-loop insulin delivery systems---so-called ``artificial pancreas'' devices---combine continuous glucose monitoring with algorithmic control of insulin delivery, further improving glycemic control and reducing the burden of diabetes management. These technological advances are direct descendants of the fundamental discovery that insulin regulates blood glucose levels.

Beyond its therapeutic use in diabetes, insulin has become one of the major tools in cell biology and biochemistry. Insulin signaling regulates a vast network of metabolic pathways, including glucose uptake (via the GLUT4 glucose transporter), glycogen synthesis, lipogenesis, and protein synthesis. The insulin receptor, a tyrosine kinase receptor, was one of the first growth factor receptors to be cloned and characterized, and the signaling pathways it activates---including the PI3K-Akt pathway and the Ras-MAPK pathway---are central to understanding cell growth, differentiation, and cancer. Dysregulation of insulin signaling is implicated in obesity, metabolic syndrome, cardiovascular disease, and certain forms of cancer, making insulin biology a major focus of biomedical research.

The production of recombinant human insulin in the 1980s, using genetically engineered \emph{Escherichia coli} or yeast, eliminated the dependence on animal pancreata and opened the door to the development of insulin analogs with modified pharmacokinetic properties. More recently, research into the regeneration of pancreatic beta cells and the transplantation of islets of Langerhans has offered hope that type 1 diabetes may one day be cured rather than merely managed. The development of stem cell-derived beta cells, which can be generated from pluripotent stem cells in the laboratory, represents a particularly promising approach that may eventually provide an unlimited supply of insulin-producing cells for transplantation.

\section{Legacy}

The legacy of Banting and Macleod is inseparable from the story of insulin. Frederick Banting, who died in a plane crash in 1941 at the age of 49, is remembered as a notable medical scientists of the twentieth century. The Banting and Best Diabetes Research Centre at the University of Toronto, the Banting Research Foundation, and the countless lives saved by insulin stand as enduring monuments to his achievement. The Canadian government has honored Banting with a place on the Canadian ten-dollar bill and with a national holiday (Banting Day), and the University of Toronto has established the Banting Postdoctoral Fellowships program to support the training of the next generation of medical researchers. Charles Best, who was denied the Nobel Prize, went on to a distinguished career in diabetes research and served as director of the Banting and Best Department of Medical Research.

John Macleod, who initially received the Nobel Prize with some reluctance on Banting's part, was a distinguished physiologist whose contributions to the discovery of insulin, while often underestimated, were essential. His laboratory, his expertise in carbohydrate metabolism, and his organizational skills provided the framework within which the discovery was made. Macleod returned to the University of Aberdeen after leaving Toronto and continued to contribute to the study of carbohydrate metabolism until his death in 1935. The insulin discovery is often cited as an example of the importance of institutional support and scientific infrastructure in enabling breakthrough research.

The discovery of insulin stands as a testament to the power of creative thinking, interdisciplinary collaboration, and perseverance in the face of uncertainty. Banting's insight---that the pancreatic duct could be ligated to destroy the exocrine tissue and liberate the islet hormone---was a notable example of ``thinking outside the box,'' and the rapid translation of the basic discovery into a life-saving therapy remains one of the fastest bench-to-bedside transitions in medical history. The legacy of insulin is not merely a single molecule but a paradigm of how basic scientific research can transform the lives of millions. Every person with diabetes who receives insulin today owes a debt of gratitude to Banting, Best, Macleod, and Collip---the four individuals whose collaborative effort produced a notable medical discoveries of all time.


\section{Modern Developments}

Insulin therapy has evolved dramatically. Insulin analogs (lispro, glargine, degludec) provide better glycemic control with fewer side effects. The artificial pancreas (closed-loop insulin delivery systems) automatically adjusts insulin based on continuous glucose monitoring, transforming diabetes management. Stem cell-derived beta cells (Vertex's VX-880) are in clinical trials for type 1 diabetes, potentially eliminating the need for insulin injections. GLP-1 receptor agonists (semaglutide, tirzepatide) have become breakthrough treatments for type 2 diabetes and obesity.
